\uuid{XdEW}
\titre{Primitive}
\auteur{bodin ; croizat ; sacre}
\organisation{exo7}
\licence{CC BY-NC-SA 4.0}
\datecreate{2022-03-08}
\niveau{L1}
\module{Analyse}
\chapitre{Calcul d'intégrales}
\sousChapitre{Primitives diverses}
\difficulte{2}
\theme{Primitive}
\exo7qcmid{hilisit.equations-differentielles.008}

\contenu{
\begin{question}
\qclassification{127.11}
Quelles sont les affirmations vraies ?
\begin{answers}
  \bad{Une primitive de $x^k$ est $\frac{x^k}{k}$.}
  \bad{Une primitive de $\ln(x)$ est $\frac{1}{x}$.}
  \good{Une primitive de $\frac{1}{\sqrt x}$ est $2\sqrt{x}$.}
  \bad{Une primitive de $e^{ax}$ est $e^{ax}$ (où $a>0$ est une constante).}
\end{answers}
\begin{explanations}
Une primitive de $x^k$ est $\frac{x^{k+1}}{k+1}$.
C'est $\ln(x)$ qui est une primitive de $\frac{1}{x}$, l'inverse est faux.
Oui, une primitive de $\frac{1}{\sqrt x}$ est $2\sqrt{x}$ puisque $(2 \sqrt{x})' = 2 \times \frac{1}{2\sqrt{x}} = \frac{1}{\sqrt{x}}$.
Enfin, une primitive de $e^{ax}$ est $\frac1a e^{ax}$.
\end{explanations}
\end{question}
}
